\scanpage{7}
On a un hom. de groupes
\[
K_s(S)\xrightarrow{\alpha}WG(S)
\]
en associant à la classe du \deleted{fibré} Module loc. libre $E$ \inserted{de r.f.}
le \deleted{\unsure{fib}} module quadratique $E\oplus\check E$, avec la forme
$\unsure{q(a,a')=\langle a,a'\rangle}$ ; ici $K_s(S)$ désigne le groupe $K$
``spécial'', groupe enveloppant du monoïde des
classes d'isomorphie de modules loc. lib. de r.f. sur $S$;
il n'est pas clair que l'hom. précédent se factorise
par $K(S)$ (déduit de $K_s(S)$ en tenant compte des
suites exactes courtes...)

(\unsure{donc} si $E$ est ext. \inserted{non splittée} de $E'$ par $E''$, il n'est pas clair
si on a $\varphi(E)=\varphi(E')+\varphi(E'')$ dans $WG(S)$....)

Par définition, $W(S)$ (groupe de Witt) est le conoyau
de l'hom. précédent
\[
K_s(S)\xrightarrow{\alpha}WG(S)\longrightarrow W(S)\longrightarrow0
\]

\uline{NB} Si $2$ inv. dans $\uline{O}_S$, alors pour tt \unsure{Module} quadratique
$(M,q)$, la somme $F=(M,q)\oplus(M,-q)$ définit \unsure{cet} \unsure{élément}
nul de $W(S)$, i.e. est $\simeq$ à un $E\oplus\check E$ comme
dessus : il suffit de prendre $E=M$, immergé
dans $F$ diagonalement et anti diagonalement
[en fait, les \inserted{mod. quad.} $F$ de la forme $(M,q)\oplus(M,-q)$ sont,
à isom. près, \deleted{\unsure{les}} de la forme $E\oplus\check E$...,
mais où $E$ est un Mod. loc. lib. de r.f.
qui admet (au moins une) forme quadratique
(str.) non dégénérée.

\unsure{De =}, si $S$ est un schém. affine, il

\mspage{6}
